\documentclass[prep,both]{debate}
\motion{辩题：「勇敢的人先享受世界」是一个陷阱 / 不是一个陷阱}
\meta{赛制}{中国科学技术大学新生辩论赛}
\meta{生成日期}{2026-09-17}
\begin{document}
\debatetitle

\begin{summary}
\begin{fields}
\field{持方优劣评估}{整体略偏正方。正方在评委直觉与材料叙事上占优：「勇敢需要条件」这个反思本身已经是流行讨论的一部分，一个没有兜底的人裸辞之后的故事，不需要论证就能打动人。反方的劣势集中在评委直觉与论证难度：它必须在评委已经觉得「这话确实有毒」的气氛里说「它不是陷阱」，而且极易滑向替这句话辩护。反方唯一稳的打法是判准之争：承认这句话不严谨，只否认它满足「陷阱」的构成要件，把战场收窄到「伪装」与「归因」两格，全程不替这句话说一句好话。}
\thesis{正方}{它把「有没有兜底」这张门票，标成了「够不够勇敢」这枚勋章，于是最付不起代价的人被叫去付账。}
\thesis{反方}{它给的是一个顺序，不是一张支票；说好的「更早开始」兑现了，没兑现的那部分是准备的事，不是这九个字的事。}
\end{fields}
\end{summary}

\section{一、赛制要点}
\begin{flowtable}
\block{A}
\flow{1 正一开篇立论}{180 秒}{正一}{二十秒内把「陷阱」的三个要件说完，否则反四质询直接冲定义}
\flow{2 反四质询正一}{120 秒 双边}{反四 问}{拿到「陷阱要有伪装」「受损必须是常态」两个承认}
\flow{3 反一开篇立论}{180 秒}{反一}{开头一句钉住质询成果，再把判准拉回「系统性受损」}
\flow{4 正四质询反一}{120 秒 双边}{正四 问}{逼反方给出「这句话到底说了什么」的正面回答}
\block{B}
\flow{6 反二驳论}{120 秒}{反二}{打正方最难圆的一步：从「不准确」到「是陷阱」中间那一跳}
\flow{7 正二驳论}{120 秒}{正二}{打反方两论点的拆台处，预留临场位回应反二}
\flow{8 二辩对辩}{各 90 秒 单边}{正二 反二}{比谁的问题对方没答，不比嗓门}
\block{C}
\flow{10 正三盘问反二反四}{90 秒 单边}{正三 问}{检验反方两人对「陷阱要不要设局者」是否同一个答案}
\flow{11 反三盘问正二正四}{90 秒 单边}{反三 问}{检验正方两人对「受损是不是常态」给不给得出证据}
\flow{12 正三盘问小结}{90 秒}{正三}{把盘问成果接到「门票换勋章」上，并预压反三}
\flow{13 反三盘问小结}{90 秒}{反三}{可直接回应正三小结，预留约 80 字临场位}
\block{D}
\flow{15 自由辩论}{各 240 秒}{全员}{主战场推进到位，必问问完，不被对方开的新战场牵走}
\block{E}
\flow{17 反四结辩}{210 秒}{反四}{封路式：提前点破正四最可能讲的那个故事}
\flow{18 正四结辩}{210 秒}{正四}{留 30–40 秒真回应反四，不背稿}
\end{flowtable}
质询是双边计时，被质询方拖时间就是在吃对方的问数，但评委看得出来，一辩答题要短而不躲。对辩是单边计时，谁都拖不了谁，比的是每次发言的密度。反四先结辩、正四后结辩，这个顺序决定了反方要写封路式，正方要留临场位。

\section{二、辩题性质与争议焦点}
\begin{fields}
\claim[辩题性质]{事实判断，带价值倡导的尾巴。句式是「X 是 / 不是 Y」，真正争的是「陷阱」这个概念的构成要件在这句话上齐不齐。}
\field{正方倾向把题目解释成}{价值倡导题：这句话在社会上起了什么作用、值不值得跟。这样正方可以用「它对谁不公平」直接调动评委。}
\field{反方倾向解释成}{严格的事实判断题：一个概念要不要适用，要看要件齐不齐；不齐就是不齐，哪怕这句话确实讨人厌。}
\field{争议焦点}{① 「陷阱」要不要设局者；② 照做的人受损是不是常态；③ 受损该算在「这句话」头上，还是算在「鲁莽」「没准备」头上。}
\field{辩论对象}{是这句 2023 年起流行的网络流行语本身，不是「勇敢」这种品质。双方都不要把题目打成「该不该勇敢」。}
\end{fields}

\section{三、关键词定义}
\subsection{3.1 陷阱（全场胜负词）}
\begin{fields}
\field{调研}{\sub{词典义}{① 为捉捕野兽而掘的深坑；② 害人的计谋。简明释义「比喻使人受骗上当的圈套」。《淮南子·兵略》「虎豹不动，不入陷阱」；《汉书·食货志下》「夫县法以诱民，使入陷阱」。（汉典《国语辞典》「陷阱」词条）}\sub{扩展用法}{现代汉语里大量存在没有设局者的「陷阱」：中等收入陷阱、贫困陷阱、消费陷阱、思维陷阱。中等收入陷阱由世界银行经济学家 Indermit Gill 与 Homi Kharas 在 2007 年报告《An East Asian Renaissance》中提出，没有任何人设这个局。}\sub{共同内核}{两种用法都要求同一件事：表面是通路，进去才发现是坑，而且是它的外观让你走进去的。}}
\begin{procard}{正方有利定义}\definition{陷阱：任何「呈现为好处、照做却让人系统性受损」的机制，不要求有设局者。}认知陷阱、结构陷阱都算。\sub{有利之处}{把「有没有人故意骗你」这一要件拿掉，正方只需证明这句话在事实层误导、照做的人受损。}\sub{反方如何反驳它}{这样一来任何不够严谨的劝告都成了陷阱。「早睡早起身体好」对倒班工人也不成立，它是陷阱吗？定义一旦失去区分力，「是陷阱」就不再是一个可被证伪的判断，等于把辩题问成了「这句话有没有缺点」。}\end{procard}
\begin{concard}{反方有利定义}\definition{陷阱：圈套。必须有设局者、有欺骗意图、有被骗上当的受害者，三者齐备。}\sub{有利之处}{这句流行语无主、无对价、无承诺，三个要件一个都不满足，反方一句话得胜。}\sub{正方如何反驳它}{这是定义杀。汉语里「中等收入陷阱」「消费陷阱」「思维陷阱」都没有设局者，按这条定义它们全都不叫陷阱，不符合语言事实；而且它让正方持方从一开始就不可能成立，评委会认为反方在躲辩论。}\end{concard}
\begin{neutralcard}{中立定义}陷阱：它呈现出来的样子和照做之后的真实后果系统性地不一致，而且正是这层外观把人引了进去。有没有具体的设局者不影响它是不是陷阱，但「受损」必须是常态，不能是个别人运气不好。\sub{合理性}{① 保留了词典义的内核（伪装加受害），同时不排除「消费陷阱」这种汉语里公认的无主体用法；② 两方都有得打，正方打「外观与真实因果不一致 + 受损是常态」，反方打「没有那层伪装」或「受损不能算在这句话头上」；③ 贴近常识，评委一听就点头。}\end{neutralcard}
\end{fields}

\subsection{3.2 勇敢的人}
\begin{fields}
\field{调研}{\sub{词典义}{勇敢：有勇气，有胆量。勇气针对的是已知的风险；不知道有风险而往前冲，汉语叫莽撞。}}
\begin{procard}{正方有利定义}\definition{勇敢的人：这句话所召唤的那种人，也就是听到它就去裸辞、去说走就走的普通年轻人。}\sub{有利之处}{把讨论落到真实听众身上，受损才有承载者。}\sub{反方如何反驳它}{这是把「鲁莽的人」换成了「勇敢的人」，偷换了主语。}\end{procard}
\begin{concard}{反方有利定义}\definition{勇敢的人：知道后果、算过账，仍然选择行动的人；没准备好就上的一律是鲁莽，不在此列。}\sub{有利之处}{一切损失都被划到「鲁莽」名下，归因这一格反方全拿。}\sub{正方如何反驳它}{那这句话就成了「准备好的人先享受世界」，等于当场承认门槛存在，正方论点一直接成立。}\end{concard}
\begin{neutralcard}{中立定义}勇敢的人：在后果不确定时仍然选择行动的人。两条边界要说明：勇敢不等于鲁莽（鲁莽是不计后果，勇敢是知道后果仍然选），勇敢也不等于有钱（这恰恰是全场最大的争点）。\sub{合理性}{同时保留了词典义与辩题的现实所指，不预先判给任何一方。}\end{neutralcard}
\end{fields}

\subsection{3.3 先享受世界}
\begin{fields}
\field{调研}{\sub{先}{顺序词，比不勇敢的人更早，不是「唯一」，也不是「更多」。}\sub{享受世界}{得到世界能给的好处与体验，旅行、恋爱、机会、自由。}}
\begin{procard}{正方有利定义}\definition{先享受世界：一种竞争性的时间承诺，不现在做就轮不到你。}\sub{有利之处}{坐实「诱导」这一要件。}\sub{反方如何反驳它}{「先」只是比较顺序，原句没有说名额有限，这是正方替它加戏。}\end{procard}
\begin{concard}{反方有利定义}\definition{先享受世界：更早开始体验，仅此而已，不含任何回报保证。}\sub{有利之处}{承诺被压到最小，落差无从谈起。}\sub{正方如何反驳它}{那这句话就什么都没说，可它能让人辞职，也被骂了三年。一句什么都没说的话不会有这种力量。}\end{concard}
\begin{neutralcard}{中立定义}先享受世界：比不敢行动的人更早得到世界给的好处。\sub{合理性}{贴住原句字面，既不加稀缺，也不把「享受」抹成零。}\end{neutralcard}
\end{fields}

\subsection{3.4 隐含要素}
\begin{fields}
\field{主体}{听到这句话并可能照做的普通年轻人，也就是没有稳定兜底、正在决定要不要裸辞、要不要 gap 的人。双方都不该拿富二代或极端个例当主体。}
\field{范围}{当下中国的网络语境。这句话 2023 年国庆假期起作为旅游九宫格配文流行，随后泛化成形容各种「说走就走」的梗。}
\field{时间尺度}{长期。短期看「先享受」大多成立，人确实去玩了；争议在几年之后的那笔账。}
\field{比较对象}{与「没有这句话、自己按条件做决定」的同一个人比，不是与「永远不行动的人」比。}
\end{fields}

\begin{warnbox}{3.5 定义杀陷阱：双方都要背}
\begin{fields}
\begin{procard}{正方的定义杀}把「陷阱」泛化成「有缺点的说法」。一旦评委听出「按你方标准什么都是陷阱」，三个要件同时塌掉。正方必须自己守住「受损是常态」这条线，并主动给出「省略只是不严谨，顶替才是陷阱」的区分。\end{procard}
\begin{concard}{反方的定义杀}把「陷阱」窄化成「诈骗」，要求正方指出设局者。评委会认为反方在用定义取消辩论。反方必须在一辩稿里主动承认「陷阱可以没有设局者」，然后在「伪装」和「系统性受损」两格上赢。\end{concard}
\end{fields}
\end{warnbox}

\section{四、判准与论证义务}
\begin{fields}
\claim[判准是什么]{要把一句话叫做陷阱，得同时看三件事：一、它呈现的因果关系和真实的因果关系系统性地不一致；二、相信它、照着做的人，受损是常态而不是例外；三、这份受损正是被它遮住的那部分东西造成的。}
\begin{procard}{正方有利判准}\definition{只要一个说法把「能承担代价」这个必要条件换成了「勇气」这个品质，并遮住了真实门槛，它就是陷阱，不必统计有多少人真的摔了。}\sub{有利之处}{免掉「系统性受损」的举证责任，正方只需在因果结构上做文章。}\sub{反方如何反驳它}{这把「不严谨」直接等同于「陷阱」。按这条标准，「努力就有回报」「越努力越幸运」全是陷阱。一个让所有励志话都变成陷阱的标准，衡量的不是陷阱，是「有没有例外」。}\end{procard}
\begin{concard}{反方有利判准}\definition{必须指认出设局者、欺骗意图与具体受害人群，三样齐了才叫陷阱。}\sub{有利之处}{这句流行语三样全缺，反方一句话得胜。}\sub{正方如何反驳它}{这是把「陷阱」改写成了「诈骗」。中等收入陷阱、消费陷阱都没有设局者，按这条标准它们全不成立；而且在这条标准下正方不可能赢，评委会认为反方用定义取消了比赛。}\end{concard}
\begin{neutralcard}{中立判准}三件同时成立，它就是陷阱；缺一件，它就只是一句不够严谨的话。\sub{合理性}{① 这是「陷阱」在汉语两种用法里共有的内核，不偏向任何一方；② 它区分了「说得不准确」和「是个坑」，正是本题的真问题；③ 双方各有明确的活干，不靠谁抢定义。}\sub{正方需要证明}{三件都成立。重点在第二与第三件：证明这句话误导人还不够，还要证明照做的人普遍因此受损，并且受损来自被它遮住的东西。}\sub{反方需要证明}{打掉任意一件即可。最省力的是第一件（它没有承诺任何回报，何来伪装）或第三件（受损来自鲁莽与无准备）。}\end{neutralcard}
\field{双方共同的红线}{不要滑向「这句话好不好」。正方不必证明它有害无益，只需证明它是个坑；反方不必替它辩护有多对，只需证明它不是坑。}
\end{fields}

\prosection{五、正方论点}
\prosubsection{论点一：门票换勋章}
\begin{fields}
\claim{这句话把「有没有兜底」这张门票，标成了「够不够勇敢」这枚勋章。}
\begin{mechanism}
\step 能不能说走就走，本来是一道算得出来的题：存款够几个月、家里要不要你补贴、简历断一年补不补得回来。
\step 这句话把这道题换成了「你够不够勇敢」，一道没有答案、只能靠自我感觉回答的题。
\step 换题之后，付不起的人不再去算价钱，只去比胆量。
\end{mechanism}
\begin{evidence}{学理}[有把握]基本归因错误（Lee Ross，1977）\sub{核心主张}{人系统性地高估个人特质的作用，低估处境的作用。}\sub{支撑}{机制第二步：把「有没有条件」这个处境变量，换成了「勇不勇敢」这个特质变量。}\sub{边界}{Ross 描述的是观察者对他人的归因偏差；用在自我归因上，要说清是「照着这套叙事看自己」，不要说成「人一定会这样想自己」。}\sub{出处}{Ross, L. (1977). The intuitive psychologist and his shortcomings. \textit{Advances in Experimental Social Psychology}, 10, 173–220.}\end{evidence}
\begin{evidence}{数据}[已核实]23.8\% 的职场人表示不会 GAP 三个月以上，理由是经济压力；另有 18.5\% 表示不会，担心简历断档增加找工作难度。\sub{来源}{智联招聘《2025 春季职场人跳槽情况调研报告》，2025 年 3 月发布，37980 份有效样本。}\sub{方法}{面向智联招聘平台用户的问卷调研，测的是跳槽意愿与离职方式，不是离职后的结果。}\sub{局限}{平台用户样本，偏向有求职行为的人群；本次通过转载全文核实，未打开智联原始 PDF。这份数据只能说明「门槛存在且人们看得见」，不能说明「谁摔了」。}\end{evidence}
\begin{evidence}{案例}[已核实]澎湃新闻 2023 年 10 月 17 日转载 Vista 看天下叶橙子的文章，记录了这句话的语义漂移：「勇敢」从「说走就走的冲动」，变成了「有说走就走的物质实力」。\sub{代表性}{记录的是这句话在传播中的用法变化，不是某一个人的遭遇，正好对应「换题」这一步。}\sub{出处}{澎湃新闻，2023-10-17，《「勇敢的人先享受世界」，普通人配吗？》。}\end{evidence}
\closing{这就是第一件事：它呈现的因果是「勇气通往先享受」，真实的因果是「兜底通往先享受」，两者系统性不一致。}
\begin{clash}
\pair{你方自己的数据证明人照样在算账，四成多人算完就不走了，伪装在哪里？}{我方从不主张每个人都不再算账。我方主张的是这句话把一道能算的题换成了一道不能算的题。而算完不走的那四成多人，正是被这句话判为「不勇敢」的人。请对方回答：一个把「我算过账，我走不了」翻译成「你不够勇敢」的说法，对这四成人是什么？}
\pair{按你方标准，所有格言都是陷阱，格言的定义就是省略条件。}{不是所有省略都是陷阱。省略了决定性变量、并且拿另一个变量顶上去，才是。「早睡早起身体好」没有告诉你睡不好是因为你不自律；这句话告诉你，你没先享受是因为你不够勇敢。省略是不严谨，顶替才是陷阱。}
\end{clash}
\end{fields}

\prosubsection{论点二：删过的账本}
\begin{fields}
\claim{它给的是一份删掉了失败样本、也删掉了反面证据的账本。}
\begin{mechanism}
\step 勇敢之后摔了的人不发帖，于是这句话的证据基础天然只剩成功样本。
\step 它靠「你以后会后悔没去」这个共识说服人，而这个共识在最大的公开复制里没有复制出来。
\step 照着一份被删过两次的账本估计自己的胜率，必然高估。
\end{mechanism}
\begin{evidence}{学理}[已核实]幸存者偏差（Abraham Wald，1943）\sub{核心主张}{只看通过了筛选的样本，会把结论推反。二战时军方只能看到飞回来的飞机，Wald 指出该加装甲的恰恰是这些飞机身上没有弹孔的地方。}\sub{支撑}{机制第一步：我们听到的「我勇敢了，我过得很好」，全部来自飞回来的那一架。}\sub{边界}{它说明的是「你的估计有偏」，不等于「真实胜率一定很低」。正方不要把它说成「勇敢一定会失败」。}\sub{出处}{Wald, A. (1943). \textit{A Method of Estimating Plane Vulnerability Based on Damage of Survivors}. Statistical Research Group, Columbia University；该备忘录 1981 年由 Center for Naval Analyses 公开。}\end{evidence}
\begin{evidence}{数据}[已核实]在迄今最大的一次公开复制里，长期回顾中 51\% 的人更困扰于自己做过的事，49\% 更困扰于自己没做的事。1994 年那条「长期看更后悔没做」的著名结论没有复制出来。\sub{来源}{Richardson, J. \& Gilovich, T. (2023). A very public replication of the temporal pattern to people's regrets. \textit{Royal Society Open Science}. DOI 10.1098/rsos.221574，n = 2600。}\sub{方法}{在公开展览现场收集参与者手写的遗憾卡片，分别问「过去几天」与「一生」两个时间窗，比较行动遗憾与不行动遗憾。}\sub{局限}{现场自选样本，非随机抽样；它证明的是「长期差异消失」，不是「做了更后悔」。正方不能把它说成后者。}\end{evidence}
\begin{evidence}{案例}[已核实]空窗期在今天的就业市场上不是中性的：国家统计局口径，不含在校生的 16–24 岁城镇调查失业率，2026 年 6 月为 14.9\%，仍高于上年同期。\sub{代表性}{它给出的是「断一年」这件事在当下的真实价钱，对应机制第三步的受损落点。}\sub{出处}{国家统计局 2026 年 7 月 20 日发布，财新网 2026-07-20 报道。}\end{evidence}
\closing{第二件与第三件：受损不是个别人手气差，是样本被筛过、反证被删掉之后必然发生的系统性高估，而这份高估的代价对没有兜底的人是不可逆的。}
\begin{clash}
\pair{幸存者偏差是所有成功叙事的共同属性，学长的经验帖、课本上的科学家故事，按你方标准全是陷阱。}{区别在于说不说条件。校友分享会会说「我当时家里能兜」，这句话不说。说了条件的叫经验，删了条件的才叫陷阱。}
\pair{五成一对四成九，这不正说明做和不做一样吗？一个五五开的选择怎么会是坑？}{五五开的是后悔，不是代价。请对方注意，我方引这份数据不是为了证明「做了更亏」，是为了证明这句话最主要的说服理由——你以后会后悔没去——在最好的证据下并不成立。它卖的不是亏本生意，是一张写着「稳赚」的账单。}
\end{clash}
\end{fields}

\subsection{（被淘汰的正方候选论点，交给反驳库与自由辩备用）}
\begin{enumerate}[leftmargin=1.8em]
\item 被商业收编：有真实获利方，但中立判准不要求设局者，这条赢不到判准分，且容易被反方反打成「你方在打广告不是打这句话」。
\item 「先」字制造稀缺与时间压力：属于替原句加戏，反方一句「先只是顺序」就打掉。
\item 享受被窄化成消费与打卡：是社会批评而非本题判断，容易跑题。
\item 不可证伪（做了没享受到就说你还不够勇敢）：机制与论点一重合，且攻击的是一种语气而不是一个文本。
\end{enumerate}

\consection{六、反方论点}
\consubsection{论点一：顺序不是支票}
\begin{fields}
\claim{它给的是一个顺序，不是一张支票；说好的那件事兑现了，没有落差。}
\begin{mechanism}
\step 陷阱要有一层伪装，伪装要有一个被盖住的落差，也就是「说好的」和「实际给的」之间的差。
\step 这句话说好的只有一件事：在先动和先等之间，先动的人更早开始。
\step 这件事本身是真的。没有落差，就没有伪装。
\end{mechanism}
\begin{evidence}{学理}[已核实]汉语「陷阱」的构成要件\sub{核心主张}{词典义是「比喻使人受骗上当的圈套」；即便是没有设局者的用法，例如世界银行 Gill 与 Kharas 2007 年提出的「中等收入陷阱」，也指一个进去就出不来的稳定状态。叫陷阱，要么有骗，要么有出不来的坑底。}\sub{支撑}{机制第一步：把「陷阱」的最低要求说清楚，正方必须交出那个被盖住的东西。}\sub{边界}{不要走到「必须有设局者」那一步，那是定义杀，评委会认为反方在躲。反方在一辩稿里就主动承认陷阱可以没有设局者。}\sub{出处}{汉典《国语辞典》「陷阱」词条；Gill, I. \& Kharas, H. (2007). \textit{An East Asian Renaissance: Ideas for Economic Growth}. World Bank.}\end{evidence}
\begin{evidence}{数据}[已核实]34.8\% 有跳槽意愿的职场人会选择裸辞，其中 21.2\% 的理由是「休息后再出发」；被问到会不会 GAP 三个月以上时，21.8\% 表示会。\sub{来源}{智联招聘《2025 春季职场人跳槽情况调研报告》，2025 年 3 月发布，37980 份有效样本。}\sub{方法}{面向智联招聘平台用户的问卷调研，测的是离职方式与理由。}\sub{局限}{它证明的是这些人各有理由、心里有数，不证明他们后来过得好；反方不许把它说成「他们都过得很好」。}\end{evidence}
\begin{evidence}{案例}[已核实]这句话最初是 2023 年国庆假期旅游九宫格的配文（「远离内耗的人和事，勇敢的人先享受世界」），随后才泛化成形容各种说走就走的梗。\sub{代表性}{说明它的原生场景是表达而不是交易：没有对价，没有签名，没有收款方。}\sub{出处}{澎湃新闻，2023-10-17，《「勇敢的人先享受世界」，普通人配吗？》。}\end{evidence}
\closing{第一件事不成立：没有那层被盖住的落差，就没有伪装，也就没有陷阱。}
\begin{clash}
\pair{「更早开始的人确实更早开始了」是同义反复，等于你方承认它兑现的是零。}{兑现的不是零，是「开始」本身，而开始是有价值的实体。请对方注意举证责任的方向：中立判准要求的是「呈现的样子与真实后果系统性不一致」，这是正方的活。请正方回答，这句话盖住的那个东西，具体是哪几个字？}
\pair{一句什么都没说的话，为什么能让人辞职，又被骂三年？}{我方从不说它什么都没说。它说了一件事：别把一生都花在等条件齐了。这句话可以被反对，但反对一句话和识破一个圈套，是两回事。}
\end{clash}
\end{fields}

\consubsection{论点二：账记在准备上}
\begin{fields}
\claim{让人摔的是没留退路，不是有勇气；把鲁莽的账记给勇敢，才是这场辩论里真正的偷换。}
\begin{mechanism}
\step 受损来自没兜底还不留退路，那是鲁莽。
\step 勇敢的词典义是知道后果仍然选，勇气针对的是已知的风险。
\step 中立判准第三件要求受损由这句话遮住的东西造成，而这九个字里没有一个字让你别算账。
\end{mechanism}
\begin{evidence}{学理}[已核实]自我效能理论（Albert Bandura，1977）\sub{核心主张}{能力信念最可靠的来源是亲历的成功经验。原文：「This source of efficacy information is especially influential because it is based on personal mastery experiences.」}\sub{支撑}{机制第一步的反面：行动本身产出的是真东西，不是空头支票。}\sub{边界}{Bandura 同样指出反复失败会降低效能感，早期就摔尤其伤。所以反方主张的是「行动有真收益」，不是「行动必然更好」。}\sub{出处}{Bandura, A. (1977). Self-efficacy: toward a unifying theory of behavioral change. \textit{Psychological Review}, 84, 191–215.}\end{evidence}
\begin{evidence}{数据}[已核实]1994 年 Gilovich 与 Medvec 收集的 213 项遗憾中，「没做的事」比「做过的事」多出接近两比一，63\% 对 37\%；另一项研究里，回顾一生时 84\% 的人对没做的事更后悔。2023 年近 2600 人的公开复制中，长期那一半变成了 51\% 对 49\%。\sub{来源}{Gilovich, T. \& Medvec, V. H. (1994). The temporal pattern to the experience of regret. \textit{Journal of Personality and Social Psychology}, 67(3), 357–365；Richardson \& Gilovich (2023), DOI 10.1098/rsos.221574。}\sub{方法}{1994 年为电话调查、书面问卷与面对面访谈，样本包括学生、职员、荣休教授与养老院居民；2023 年为公开展览现场收集的手写卡片。}\sub{局限}{反方只能主张「做和不做，长期看后悔的机会一样大」，不能只报 1994 不报 2023，否则对方当场可以打伪证。}\end{evidence}
\begin{evidence}{案例}[已核实]同一句话，两种结局：被骂的版本几乎都是「裸辞去大理然后没钱了」，没有人骂「攒够两年生活费再辞职去读书」。智联的同一份报告里，裸辞的人中 21.2\% 说的是「休息后再出发」，这正是算过账的那一类。\sub{代表性}{对比的是两种准备程度，而不是两种胆量，正好落在归因这一格。}\sub{出处}{智联招聘《2025 春季职场人跳槽情况调研报告》。}\end{evidence}
\closing{第三件事不成立：受损的归因落在准备上，不落在这九个字上。}
\begin{clash}
\pair{这句话里哪个字提醒你去算后果？你方是在替它加注脚。}{注脚不是我方加的，是词典写的。勇敢的意思就是有勇气有胆量，而勇气针对的是已知的风险；不知道有风险还往前冲，汉语里叫莽撞。对方要打的那种人，题目里没有。}
\pair{五五开只是后悔打平，成本没减：做的那一边要掏钱、掏时间、掏简历上的一年。}{后悔研究测的正是事后的总体评价，一个人回头看，把花掉的钱和断掉的一年全算进去之后，还是这么答的。成本已经在里面，不能再减一次。而且对方这一步等于亲口承认：胜负在有没有算账。那就回到我方第二点——算账是准备的事。}
\end{clash}
\end{fields}

\subsection{（被淘汰的反方候选论点，交给反驳库与自由辩备用）}
\begin{enumerate}[leftmargin=1.8em]
\item 「没有统一的坑」：举证责任本就在正方，主动扛下来是自找麻烦，留作自由辩的追击句。
\item 「叫它陷阱的代价」：论证的是「不该这么叫」而不是「它不是」，作为价值层留给反四结辩。
\item 「每句励志话都有例外」：单独说像狡辩，只作反驳口径，不作论点。
\item 「它只是配图文案，没人真照它辞职」：会被「那它为什么红三年」反打，降格为论点一的支撑。
\end{enumerate}

\prosection{七、正方反驳库（针对反方全部可能论点）}
\begin{rebuttals}
\rebuttal{1}{它给的是一个顺序，不是一张支票}{顺序也是承诺，它承诺的是「先」。}{打机制。对方说这句话兑现了「更早开始」，可「更早开始」是同义反复，等于零兑现。真正让人听进去的是「享受」两个字，不是「开始」两个字。请评委听一遍原句，是「勇敢的人先享受世界」，不是「勇敢的人先开始」。对方替它改了半句。}{对方问「盖住的那个东西是哪几个字」→ 答：兜底能力。这句话的因果里没有它，现实的因果里少了它谁也走不了。}{最终}
\rebuttal{2}{让人摔的是没留退路，不是有勇气}{那对方守的是「准备好的人先享受世界」，不是今天这道题。}{打定义。对方把「勇敢」注解成「知道后果仍然选」，可这九个字里没有一个字提到后果。如果这句话真的包含「先算账」，请对方指出是哪个字。指不出来，那条注脚就是对方替它加的，而加注脚本身正好说明原句缺了它。}{对方问「难道词典义不算」→ 答：算，词典义正是我方的论据。勇气针对已知风险，而这句话不提供任何风险信息，它只提供一个称号。}{最终}
\rebuttal{3}{陷阱必须有设局者，这句话无主无对价}{中等收入陷阱、消费陷阱、思维陷阱，哪一个有设局者？}{打定义。世界银行的 Gill 与 Kharas 在 2007 年提出中等收入陷阱，没有任何人设这个局。把「设局者」写进定义，等于把汉语里一半的「陷阱」判成用词错误，也等于让我方持方从第一分钟起就不可能成立。评委会看出这是定义杀。}{对方问「那按你方标准什么不是陷阱」→ 答：说了条件的经验分享不是；只省略、没顶替的格言不是。我方的线很清楚。}{常见}
\rebuttal{4}{叫它陷阱，会让本来就不敢的人更不敢}{这说的是后果，不是性质；题目问的是「是不是」。}{打判准。就算指出陷阱真的会让人却步，陷阱也还是陷阱。而且对方这一条恰恰承认了那里有个坑，否则何必担心有人不敢走。}{对方追问「那你方是在劝人别勇敢吗」→ 答：不是。我方劝的是先看清门票价格再决定买不买，这与勇敢并不冲突。}{常见}
\rebuttal{5}{每句励志话都有例外，有例外不等于是坑}{我方从不打例外，我方打的是顶替。}{省略条件只是不严谨；拿另一个变量顶替被省略的那个，才是陷阱。「早睡早起身体好」没告诉你睡不好是因为你不自律；这句话告诉你，你没先享受是因为你不够勇敢。}{对方追问「顶替有什么证据」→ 答：这句话的语义漂移本身就是证据，「勇敢」在传播中被重新定义成了「有说走就走的物质实力」。}{淘汰}
\rebuttal{6}{几千万人各有各的出口，不存在一个统一的坑}{陷阱不需要同一个坑底，只需要同一层伪装。}{打定义。捕兽的坑可以挖在不同的地方，让它们都成为陷阱的是上面那层浮土。这句话就是那层浮土：把门槛叫成勇气。}{对方追问「那你方怎么证明受损是常态」→ 答：受损的常态就是代价错配，付不起的人按付得起的人的活法行动，这在没有兜底的人身上是普遍的。}{淘汰}
\rebuttal{7}{1994 年的研究说长期看人更后悔没做的事}{请对方把 2023 年那份复制一起念完。}{打论据。Richardson 与 Gilovich 2023 年在 Royal Society Open Science 上发表的近 2600 人公开复制里，长期回顾是 51\% 对 49\%，那条著名结论没有复制出来。对方引的是三十年前一个小样本上被推翻的那一半。}{对方说「五五开就不是陷阱」→ 答：五五开的是后悔，不是代价。我方论点二说的正是这个：它给的是一份删掉了反面证据的账本。}{常见}
\rebuttal{8}{它只是朋友圈配图文案，没人真照它辞职}{那它为什么能红三年？}{一句真的什么都没影响的话不会流行。对方在这里二选一：要么承认它有影响力，那它就说了点什么；要么坚持它没有，那对方等于承认我们在讨论一句废话，可废话不会被骂三年。}{对方说「流行不等于有效」→ 答：那请对方解释，为什么被骂的全是「照它做了然后没钱了」的人，而不是别的什么人。}{常见}
\end{rebuttals}

\consection{八、反方反驳库（针对正方全部可能论点）}
\begin{rebuttals}
\rebuttal{1}{这句话把门票换成了勋章}{对方自己的数据证明人照样在算账。}{打论据。智联三万七千多份样本里，四成多人明说不 GAP，理由是经济压力和怕简历断档，他们一个字不差地算了账。一个让四成多人当场算出自己走不了的说法，伪装在哪里？对方说它「取消了算账这个问题」，可现实里这个问题一次都没被取消。}{对方说「这句话把他们判成了不勇敢」→ 答：判他们的不是这九个字，是对方的读法。原句只说了勇敢的人怎样，没说不勇敢的人是谁。}{最终}
\rebuttal{2}{它给的是一份删过的账本，幸存者偏差}{按这个标准，今天所有的经验分享都是陷阱。}{打定义。幸存者偏差是所有成功叙事的共同属性，校友分享会、学长的经验帖、课本上的科学家故事，无一例外。对方需要的不是「有偏差」，是「受损是常态」。对方全场没有给出一条「照这句话做的人普遍变差」的证据。}{对方说「区别在于顶替」→ 答：那请对方证明有人真的因为这句话放弃了算账。拿不出来，这一步就还是一个形容词。}{最终}
\rebuttal{3}{它被旅游和分期收编，有真实获利方}{那对方打的是广告，不是这九个字。}{打归因。有人用这句话卖东西，和这句话本身是不是陷阱，是两件事。任何流行语都会被拿去卖东西，「好好爱自己」也在卖东西，它是陷阱吗？}{对方说「获利方证明了诱导」→ 答：诱导的是广告主，不是这句话。请对方把被告席上的人换对。}{淘汰}
\rebuttal{4}{「先」字制造稀缺和时间压力}{「先」是顺序词，不是限量。}{打定义。原句没有说「只有勇敢的人能享受」，也没有说名额有限。对方替它加了一个它没有的意思，这在辩论里叫加戏。}{对方说「先本身就有紧迫感」→ 答：那「先来后到」也在制造焦虑吗？顺序词的紧迫感是对方带进来的。}{淘汰}
\rebuttal{5}{享受被窄化成旅游打卡和消费}{那是在批评消费主义，不是在判断这句话。}{打判准。这一条就算全部成立，得到的结论也是「有人误用了它」，不是「它是陷阱」。中立判准问的是三件具体的事，这一条一件都没回答。}{对方说「误用足以说明它有诱导性」→ 答：误用说明的是使用者，不是文本。按这个标准，被用坏的每一句话都是陷阱。}{淘汰}
\rebuttal{6}{不可证伪：做了没享受到就说你还不够勇敢}{会这么说的是人，不是这句话。}{打归因。原句里没有任何一处写「如果没成功就是你不够勇敢」。对方攻击的是一种语气，不是一个文本；而语气是使用者带进来的。}{对方说「这就是它的话术结构」→ 答：请对方把这个结构从九个字里指出来。指不出来，那就是对方自己补全的。}{淘汰}
\rebuttal{7}{青年失业率高，空窗期代价很高}{代价高，说明要准备，不说明它是陷阱。}{打机制。失业率高恰恰证明「有没有准备」这个变量权重很大，而准备不准备是人自己的事，不是这句话的事。也请对方把趋势念完：国家统计局口径，不含在校生的 16–24 岁失业率 2026 年 6 月是 14.9\%，较上月下降 0.7 个百分点，已经连续三个月回落，虽然仍高于上年同期。}{对方说「对没兜底的人代价不可逆」→ 答：承认。所以我方从第一分钟就说，这句话不负责替人免掉准备，它只是把顺序说了一遍。}{常见}
\rebuttal{8}{说到底这就是消费主义话术}{给一句话贴标签，不等于证明它的性质。}{打判准。中立判准问的是三件具体的事：有没有伪装、受损是不是常态、账该记给谁。「消费主义话术」这五个字一件都没回答。请对方回到判准上来。}{对方说「消费主义本身就是伪装」→ 答：那请对方指出这句话向谁收了钱、承诺了什么。指不出来，伪装这一格就是空的。}{常见}
\end{rebuttals}

\section{九、持方优劣评估}
\begin{assessment}
\assess{定义空间}{「陷阱」的词典义是「使人受骗上当的圈套」，核心带着「骗」，这一层天然靠近反方；但汉语里「中等收入陷阱」「消费陷阱」这类无主体用法同样是公认的，正方拿它就能把设局者要件拆掉。两边都有落脚点，谁先在一辩稿里把中立定义说透，谁就拿到主场。}{均衡}
\assess{论证义务}{中立判准下正方要立三件，反方打掉任意一件即可，这是结构性的不对称。正方最难的是第二件「受损是常态」：这句话流行才三年，没有任何一项研究测过「相信它的人后来怎么样了」，正方注定拿不出直接证据。}{偏反}
\assess{论据可得性}{公开讨论的材料几乎一边倒地站在正方：语义漂移的记录、经济门槛的数据、就业市场的数字都现成。反方能用的硬材料只有后悔研究与自我效能，而后悔研究最有名的那一半恰好在 2023 年的复制里没站住，反方必须主动交出这一点才不至于被打伪证。}{偏正}
\assess{评委直觉}{未经论证时，普通听众更容易点头的是「这话确实有毒」，因为「勇敢需要条件」这个反思本身就是这几年流行讨论的一部分。反方要先克服这个直觉。另一方面，新生赛评委多为高年级辩手与教练，普遍反感定义杀和「什么都叫陷阱」的泛化，正方如果把口子开得太大，直觉优势会当场翻转。本项只用一般性判断，没有任何关于具体评委偏好的信息。}{略偏正}
\assess{赛制顺序}{正方先立论，能第一个把「陷阱三要件」定下来，这在定义之争里是实打实的便宜；反方先结辩，可以写封路式，提前拆掉正四最可能讲的那个故事。两项大致抵消。}{均衡}
\end{assessment}
\begin{fields}
\claim[结论]{整体略偏正方。反方较劣势，劣势主要来自评委直觉，其次来自论据可得性；反方的优势全部集中在论证义务这一格，所以反方的全部打法都应该从那一格出发。}
\begin{concard}{劣势方（反方）打法}\textbf{收窄战场}：只打两格，一是「伪装」，二是「归因」，不接「这句话好不好」「消费主义」「年轻人该不该躺平」这些战场，用一句话带回判准。\sub{抢先定义}{反一开头必须主动承认「陷阱可以没有设局者」，把定义杀的嫌疑自己摘掉，评委才会听后面的话。}\sub{判准之争}{全场反复问同一个问题：这句话盖住的那个东西，具体是哪几个字？正方答不出「被盖住的 A」，伪装这一格就是空的。}\sub{价值层}{反四结辩的升华不要替这句话辩护，要写「把一句劝人别一直等下去的话定性成陷阱，代价落在谁身上」，落到一个具体的人身上。}\sub{顺序利用}{反四先结辩，提前点破正四最可能的收束方式（讲一个没兜底的人裸辞之后的故事），让评委听正四时产生「这个反方刚才说过了」的感觉。}\end{concard}
\begin{procard}{优势方（正方）要防的}最容易翻车的地方是把口子开得太大。一旦评委听出「按正方标准什么都是陷阱」，前面所有的直觉优势会在三十秒内清零。正方必须主动、反复地说那条限定：省略是不严谨，顶替才是陷阱。第二个风险是被拖进「你方是不是在劝人别勇敢」，那不是本题，要一句话带回。\end{procard}
\end{fields}

\section{十、口径表}
\prosubsection{10.1 正方口径表}
\begin{canon}{正方}
\canongroup{定义}
\canonrow{陷阱}{它呈现出来的样子和照做之后的真实后果系统性地不一致，而且正是这层外观把人引了进去}{看着是路，走进去是坑}{只要有缺点就是陷阱；一定有人故意骗你}{汉典《国语辞典》「陷阱」词条}
\canonrow{无需设局者}{中等收入陷阱、消费陷阱都没有人设局，它们照样叫陷阱}{坑不一定有人挖}{一定有人在背后获利}{Gill \& Kharas 2007}
\canonrow{勇敢的人}{在后果不确定时仍然选择行动的人}{敢往前走的人}{不计后果的人}{《现代汉语》词典义}
\canonrow{先享受世界}{比不敢行动的人更早得到世界给的好处}{更早拿到好东西}{唯一能享受；名额有限}{原句}
\canongroup{判准}
\canonrow{三件事}{一、外观与真实因果系统性不一致；二、照做的人受损是常态；三、受损正是被它遮住的东西造成的}{伪装、受损、归因}{只要不严谨就是陷阱}{备赛文档第四节}
\canonrow{我方义务}{三件都要立，举证责任在我方}{我方扛全部举证}{对方也要举证（对方只需打掉一件）}{同上}
\canonrow{顶替线}{省略条件只是不严谨，拿另一个变量顶替被省略的那个才是陷阱}{省略不算，顶替才算}{所有励志话都是陷阱}{同上}
\canongroup{论点}
\canonrow{门票换勋章}{这句话把「有没有兜底」这张门票，标成了「够不够勇敢」这枚勋章}{把能不能，换成了敢不敢}{穷人不配勇敢}{Ross 1977}
\canonrow{删过的账本}{它给的是一份删掉了失败样本、也删掉了反面证据的账本}{只让你看飞回来的那一架飞机}{勇敢一定会失败}{Wald 1943}
\canongroup{数据}
\canonrow{门槛数据}{23.8\% 的职场人不会 GAP 三个月以上，理由是经济压力；另有 18.5\% 是担心简历断档}{稿里说「两成四」「将近两成」，合起来「四成多人」}{四成多人后悔了}{智联招聘《2025 春季职场人跳槽情况调研报告》，37980 份样本}
\canonrow{五五开}{长期回顾里 51\% 更困扰于做过的事，49\% 更困扰于没做的事}{五成一对四成九；五五开}{做了更后悔}{Richardson \& Gilovich 2023, n=2600}
\canonrow{空窗价}{不含在校生的 16–24 岁城镇调查失业率，2026 年 6 月为 14.9\%，仍高于上年同期}{一成五左右}{青年失业率创新高}{国家统计局，财新 2026-07-20}
\canongroup{案例}
\canonrow{语义漂移}{「勇敢」从「说走就走的冲动」，变成了「有说走就走的物质实力」}{勇敢改了名，门槛没变}{这句话是商家编出来的}{澎湃新闻 2023-10-17}
\end{canon}

\consubsection{10.2 反方口径表}
\begin{canon}{反方}
\canongroup{定义}
\canonrow{陷阱}{它呈现出来的样子和照做之后的真实后果系统性地不一致，而且正是这层外观把人引了进去}{看着是路，走进去是坑}{必须有设局者（定义杀，绝对不许说）}{汉典《国语辞典》「陷阱」词条}
\canonrow{承认项}{我方承认陷阱可以没有设局者}{坑不一定有人挖}{请对方指出是谁设的局}{Gill \& Kharas 2007}
\canonrow{勇敢的人}{在后果不确定时仍然选择行动的人；勇气针对的是已知的风险}{知道有风险还是走}{勇敢就是准备好了}{《现代汉语》词典义}
\canonrow{它说了什么}{它给的是一个顺序，不是一张支票：在先动和先等之间，它推荐先动}{它排了个序，没开支票}{它什么都没说}{原句}
\canongroup{判准}
\canonrow{三件事}{一、外观与真实因果系统性不一致；二、照做的人受损是常态；三、受损正是被它遮住的东西造成的}{伪装、受损、归因}{必须有欺骗意图}{备赛文档第四节}
\canonrow{我方义务}{打掉任意一件即可；三件缺一，它就只是一句不够严谨的话}{三缺一就不成立}{我方要证明这句话是对的}{同上}
\canongroup{论点}
\canonrow{顺序不是支票}{它承诺的只有「先行动的人更早开始」，这件事兑现了，没有落差}{说好的给了，哪来的坑}{它没有任何影响力}{汉典「陷阱」词条}
\canonrow{账记在准备上}{让人摔的是没留退路，不是有勇气；这九个字没有一个字让你别算账}{摔跤的是没准备，不是不怕}{摔了是活该}{Bandura 1977}
\canongroup{数据}
\canonrow{各有理由}{34.8\% 有跳槽意愿的人会选择裸辞，其中 21.2\% 是「休息后再出发」；21.8\% 表示会 GAP 三个月以上}{三成五；两成一；两成二}{他们都过得很好}{智联招聘《2025 春季职场人跳槽情况调研报告》，37980 份样本}
\canonrow{五五开}{长期回顾里 51\% 更困扰于做过的事，49\% 更困扰于没做的事}{五成一对四成九；五五开}{不做更后悔}{Richardson \& Gilovich 2023, n=2600}
\canonrow{后悔研究原版}{1994 年 213 项遗憾中，没做的事占 63\%，做过的事占 37\%；回顾一生时 84\% 的人对没做的事更后悔}{接近两比一；八成四}{只报 1994 不报 2023}{Gilovich \& Medvec 1994, JPSP 67(3), 357–365}
\canonrow{空窗价}{不含在校生的 16–24 岁城镇调查失业率，2026 年 6 月为 14.9\%，较上月下降 0.7 个百分点，连续三个月回落，仍高于上年同期}{一成五左右，还在往下走}{就业市场已经好了}{国家统计局，财新 2026-07-20}
\canongroup{案例}
\canonrow{配图文案}{这句话最初是 2023 年国庆假期旅游九宫格的配文}{一句朋友圈配文，没有签名也没有收款方}{所以它无关紧要}{澎湃新闻 2023-10-17}
\end{canon}

\section{十一、论据出处清单}
\begin{sourcelist}
\source{1}{「陷阱」的词义：捕兽的深坑；害人的计谋、使人受骗上当的圈套}{汉典《国语辞典》「陷阱」词条，\url{https://zdic.net/hant/陷阱}，2026 年访问}{比喻使人受骗上当的圈套；虎豹不动，不入陷阱}{2026-09-17}{已核实}{双方定义、双方一辩稿}{}
\source{2}{「中等收入陷阱」由世界银行经济学家 Gill 与 Kharas 于 2007 年提出，没有设局者}{英文维基百科 Middle income trap 词条，\url{https://en.wikipedia.org/wiki/Middle_income_trap}；原报告 Gill, I. \& Kharas, H. (2007). An East Asian Renaissance: Ideas for Economic Growth. World Bank}{The term was coined by economists Indermit Gill and Homi Kharas in 2007}{2026-09-17}{已核实}{正方打定义杀；反方主动承认项}{原报告 PDF 见世界银行 Open Knowledge Repository，可进一步核对原文表述}
\source{3}{这句话的语义漂移：「勇敢」从冲动变成物质实力}{澎湃新闻 2023-10-17 转载 Vista 看天下 叶橙子《「勇敢的人先享受世界」，普通人配吗？》，\url{https://m.thepaper.cn/newsDetail_forward_24969878}}{「勇敢」也逐渐从拥有说走就走的冲动，微妙地转化为有说走就走的物质实力}{2026-09-17}{已核实}{正方论点一案例}{}
\source{4}{「勇敢」需要托底：能随便辞职、gap 的人大概率有存款或家庭资助}{澎湃新闻 2023-10-17，\url{https://m.thepaper.cn/newsDetail_forward_24969878}}{能随随便便辞职、gap，说明身上没有车贷房贷之类的固定生活支出}{2026-09-17}{已核实}{正方论点一、自由辩}{}
\source{5}{这句话最初是 2023 年国庆假期旅游九宫格的配文}{澎湃新闻 2023-10-17，\url{https://m.thepaper.cn/newsDetail_forward_24969878}}{远离内耗的人和事，勇敢的人先享受世界}{2026-09-17}{已核实}{反方论点一案例}{}
\source{6}{23.8\% 不会 GAP 三个月以上（经济压力）；另有 18.5\% 担心简历断档}{智联招聘《2025 春季职场人跳槽情况调研报告》，37980 份有效样本，2025 年 3 月发布；经三茅人力资源网全文转载核实，\url{https://www.hrloo.com/news/333870.html}}{23.8\% 的职场人表示「不会」GAP，主要是出于对经济压力的担忧}{2026-09-17}{已核实}{正方论点一数据}{可向智联招聘官网索取报告 PDF 原件，核对样本构成}
\source{7}{34.8\% 有跳槽意愿者会裸辞，其中 21.2\% 是「休息后再出发」；21.8\% 会 GAP 三个月以上}{智联招聘《2025 春季职场人跳槽情况调研报告》，\url{https://www.hrloo.com/news/333870.html}}{共计 34.8\% 有跳槽意愿的职场人会选择裸辞}{2026-09-17}{已核实}{反方论点一数据、反方论点二案例}{同上}
\source{8}{2023 年公开复制：长期回顾 51\% 更困扰于做过的事，49\% 更困扰于没做的事}{Richardson, J. \& Gilovich, T. (2023). A very public replication of the temporal pattern to people's regrets. Royal Society Open Science. DOI 10.1098/rsos.221574，n = 2600}{only 51\% (n = 675) said they were more troubled by their most regrettable action}{2026-09-17}{已核实}{正方论点二数据、反方论点二数据}{}
\source{9}{1994 年原始研究：213 项遗憾中没做的事占 63\%，做过的事占 37\%}{Gilovich, T. \& Medvec, V. H. (1994). The temporal pattern to the experience of regret. Journal of Personality and Social Psychology, 67(3), 357–365. DOI 10.1037/0022-3514.67.3.357}{Regrettable failures to act outnumbered regrettable actions by nearly a 2 to 1 margin (63\% vs. 37\%)}{2026-09-17}{已核实}{反方论点二数据}{}
\source{10}{1994 年原始研究：回顾一生时 84\% 的人对没做的事更后悔}{Gilovich \& Medvec (1994), JPSP 67(3), 357–365. DOI 10.1037/0022-3514.67.3.357}{a substantial majority of the subjects (84\%) reported greater regret for what they failed to do}{2026-09-17}{已核实}{反方论点二数据、反方自由辩}{}
\source{11}{幸存者偏差：只看通过筛选的样本会把结论推反}{Wald, A. (1943). A Method of Estimating Plane Vulnerability Based on Damage of Survivors. Statistical Research Group, Columbia University；1981 年由 Center for Naval Analyses 公开；书目见 \url{https://en.wikipedia.org/wiki/Survivorship_bias}}{Wald, Abraham. (1943). A Method of Estimating Plane Vulnerability Based on Damage of Survivors}{2026-09-17}{已核实}{正方论点二学理}{备忘录全文可在 CNA 公开档案中调阅，赛前可补一句原文}
\source{12}{自我效能：能力信念最可靠的来源是亲历的成功经验}{Bandura, A. (1977). Self-efficacy: toward a unifying theory of behavioral change. Psychological Review, 84, 191–215. DOI 10.1037/0033-295X.84.2.191}{This source of efficacy information is especially influential because it is based on personal mastery experiences}{2026-09-17}{已核实}{反方论点二学理}{}
\source{13}{基本归因错误：人高估个人特质、低估处境}{Ross, L. (1977). The intuitive psychologist and his shortcomings: Distortions in the attribution process. Advances in Experimental Social Psychology, 10, 173–220}{}{}{有把握}{正方论点一学理}{本次未打开原章正文，只核到书目信息；赛前请查原章或 Ross 2018 年在 Perspectives on Psychological Science 上的回顾文，补一句原文定义}
\source{14}{不含在校生的 16–24 岁城镇调查失业率 2026 年 6 月为 14.9\%，较上月降 0.7 个百分点，连续三个月回落，仍高于上年同期}{国家统计局 2026 年 7 月 20 日发布；财新网 2026-07-20 报道，\url{https://economy.caixin.com/2026-07-20/102466293.html}}{2026 年 6 月，16—24 岁劳动力（不包含在校学生）失业率录得 14.9\%，较上月下降 0.7 个百分点}{2026-09-17}{已核实}{正方论点二案例、反方反驳库第 7 条}{可在国家统计局官网分年龄组失业率数据页复核}
\source{15}{这句话的最早出处与首发者}{网络说法不一：有指向 2024 年李梦霁同名书籍与郭文韬同名歌曲的，但两者都晚于 2023 年的流行时间，无法确认首发}{}{}{待核实}{备用，未进稿件}{查证方向：微博与小红书 2023 年上半年最早发布时间；「勇敢的人先享受世界」最早媒体报道}
\source{16}{这句话被旅游平台或消费信贷用作广告语的具体投放案例}{本次检索未找到可核实的具体投放（品牌、时间、渠道）}{}{}{待核实}{备用，未进稿件；「被商业收编」这一淘汰论点若要上场必须先补上}{查证方向：广告备案库、品牌官微历史推文、行业营销案例库}
\end{sourcelist}

\section{十二、备赛建议}
\begin{fields}
\begin{timeline}
\when{第 10–8 天}{全队背下中立定义与中立判准的三件事。两方各自明确自己主打哪两格：正方打伪装与受损，反方打伪装与归因。}
\when{第 7–5 天}{一辩稿成型并出声念三遍；二辩、四辩对齐被盘问的全部答案；把「陷阱要不要设局者」这一问的答案写成同一句话背下来。}
\when{第 4–3 天}{模辩一场，重点检验三件事：正方会不会被打成「什么都叫陷阱」，反方会不会滑成「替这句话辩护」，双方数据有没有说错。}
\when{第 2–1 天}{按模辩结果只改口径表，不换论点；压时间，每篇稿子念到比规定时间少 10 秒。}
\end{timeline}
\field{模辩重点检验}{① 正方被问「那什么不是陷阱」时，能不能三秒内说出「省略不算，顶替才算」；② 反方被问「这句话到底说了什么」时，能不能给出正面回答而不是「它什么都没说」；③ 双方引后悔研究时，有没有把 1994 年与 2023 年两组数字一起说完。}
\begin{checklist}
\item 陷阱的中立定义一句话
\item 判准的三件事，按顺序
\item 我方两个论点的标签
\item 我方那条最关键的数据，连机构带年份
\item 对方最强的一句话，以及我方的回应
\item 我方全场反复追问对方的那一个问题
\end{checklist}
\end{fields}
\end{document}
